\section{Aprendizaje basado en modelos con Dyna-Q y Prioritized Sweeping}

\subsection*{a) Aprendizaje directo y actualización del modelo}

Cada actualización de Q-learning, real o simulada, aplica la regla
\[
    Q(s,a)\leftarrow Q(s,a)
    +\alpha\bigl[r+\gamma\max_{a'}Q(s',a')-Q(s,a)\bigr].
\]
Para la transición real de $A$ a $B$ con recompensa $6.4$, el objetivo es
$6.4+0.9(4.0)=10.0$. A partir de $Q(A,a_1)=0.0$, queda
\[
    Q(A,a_1)\leftarrow 0.0+0.5(10.0-0.0)=\boxed{5.0}.
\]
El modelo guarda el resultado observado
$\operatorname{Model}(A,a_1)\leftarrow (R=6.4,S'=B)$.

\subsection*{b) Dos pasos de planificación}

En el primer paso se consulta $\operatorname{Model}(A,a_1)$ y se recupera
la misma transición.
El objetivo sigue siendo $10.0$, por lo que
\[
    Q(A,a_1)\leftarrow 5.0+0.5(10.0-5.0)=\boxed{7.5},
    \qquad Q(C,a_2)=0.0.
\]
En el segundo paso se consulta $\operatorname{Model}(C,a_2)=(0,A)$.
Se usa la tabla obtenida después del primer paso, de modo que
\[
    Q(C,a_2)\leftarrow 0.0+0.5\bigl[0+0.9\max_a Q(A,a)-0.0\bigr]
    =0.45\max_a Q(A,a).
\]
El ejercicio solo proporciona $Q(A,a_1)$. Si $a_1$ maximiza el valor en
$A$ después del primer paso, entonces
$Q(C,a_2)=0.45(7.5)=\boxed{3.375}$. Si otra acción tiene un valor mayor,
se usa ese valor en el máximo. Al terminar, $Q(A,a_1)=7.5$.

\subsection*{c) Arquitectura y selección de actualizaciones}

Dyna-Q intercala interacción y planificación. Cada transición real
actualiza directamente la tabla $Q$ y aporta un dato al modelo. A
continuación, la planificación consulta el modelo $k$ veces para generar
transiciones simuladas. Ambas fuentes de experiencia actualizan la misma
tabla con la regla de Q-learning. Los pasos simulados se realizan sin
interactuar nuevamente con el entorno~\cite[seccs.~8.1--8.3]{sutton}.

En un laberinto con recompensa solo en la meta, tras encontrarla por
primera vez, el valor empieza a propagarse hacia los pares que pueden
conducir a ella. Con muestreo uniforme se pueden gastar muchos pasos de
planificación en pares cuyo valor apenas cambiaría. Prioritized Sweeping
usa una cola de prioridades basada en la magnitud del error TD
$\lvert r+\gamma\max_{a'}Q(s',a')-Q(s,a)\rvert$ que predice el modelo.
Tras actualizar $(s,a)$, reconsidera la prioridad de los pares que
conducen a $s$. Así propaga el valor hacia atrás desde la
meta~\cite[secc.~8.4]{sutton}.
